\section{Point tests and controlled avoidance}
\label{sec:stable-avoidance}

The next step turns a stable vanishing statement into a small change of a
map.  The parameter in the test records the points to be avoided.  Keeping
this parameter throughout the argument supplies the compatibility needed
to make the changes on adjacent simplices agree.

\subsection{Point tests and their products}

We use the cubes $C(t,n)$ and $T(t,n)$ and the sparse arrangements
$Z_s(t,n)$ from \Cref{sec:introduction}, and write $M=tn$.
Coordinate subspaces below may also be specified by setting individual
scalar coordinates to zero, so the terminology includes whole-slot
conditions.  A point test records equality with a parameter $z$ as a
relative stable cohomotopy class.  Its relative support is essential
for the later control on displacement.

\begin{definition}[Point test]\label{def:point-test}
Recall $c=1/32$.  Let $P$ be a finite polyhedron, let
$f:P\to C(t,n)$ be PL, and let $Z\subseteq T(t,n)$ be a subpolyhedron.
Define
\begin{align*}
 A_c(f,Z)&=\{(x,z)\in P\times Z:\|f(x)-z\|_\infty\ge c\},\\
 Q_c(f,Z)&=(P\times Z)/A_c(f,Z).
\end{align*}
For $0<\rho<c$, take the oriented collapse
$\kappa_\rho:\R^M\to S^M$ of the complement of
$(-\rho,\rho)^M$.  The \emph{point test} is the based map
\begin{equation}\label{eq:point-test-map}
 \varphi_\rho(f,Z):Q_c(f,Z)\longrightarrow S^M,
 \qquad [x,z]\longmapsto\kappa_\rho(f(x)-z).
\end{equation}
Its stable class is denoted by
\[
 \omega_{\mathbb S}(f,Z)\in
 [\Sigma^\infty Q_c(f,Z),\Sigma^M\mathbb S].
\]
Composing with the unit $\mathbb S\to MU$ gives
$\omega_{MU}(f,Z)\in\widetilde{MU}^{M}(Q_c(f,Z))$.
\end{definition}

One concrete choice of $\kappa_\rho$ is the smash product of $M$ oriented
one-dimensional interval collapses, in the order of the scalar
coordinates.  Varying $\rho$ gives a homotopy through maps constant
outside $(-c,c)^M$.  More generally, any positively oriented collapse
supported in that cube represents its relative generator and gives the
same class.  We may therefore shrink the support whenever necessary.
The difference map is PL, so $A_c(f,Z)$ is a subpolyhedron after a finite
subdivision.  The quotient in \Cref{def:point-test} is consequently a
finite pointed polyhedron.

\begin{lemma}[Restrictions and products]\label{lem:point-test-products}
Point tests are natural under pullback in the input and under restriction
of the parameter set.  In particular, if $Z\subseteq Z'$ are
subpolyhedra of the same parameter cube, the induced map
$Q_c(f,Z)\to Q_c(f,Z')$ identifies the former with a pointed
subpolyhedron after a common subdivision.

Suppose $f_i:P_i\to[-1,1]^{M_i}$ and
$Z_i\subseteq[-2,2]^{M_i}$ are finitely many such data.  With coordinates
concatenated in order, there is a canonical based homeomorphism
\begin{equation}\label{eq:point-test-product-quotient}
 Q_c\left(\prod_i f_i,\prod_i Z_i\right)
 \cong\bigwedge_i Q_c(f_i,Z_i).
\end{equation}
Under the corresponding identification
$S^{\sum_iM_i}\cong\bigwedge_i S^{M_i}$, the point test is the smash
product of the point tests.  The analogous equality for $MU$ uses the
multiplication of $MU$.
\end{lemma}

\begin{proof}
All restriction maps are induced by maps of the defining pairs.  For the
subpolyhedron assertion, triangulate $P\times Z'$ so that $P\times Z$
and $A_c(f,Z')$ are subpolyhedra.  Collapsing the latter identifies no two
points of the former except those already collapsed in $Q_c(f,Z)$.
Quotients of this finite polyhedral pair give the assertion, including
the stated basepoint convention.

For products the relative subset is exactly the union of the subsets
on which at least one factor belongs to $A_c(f_i,Z_i)$.  Its quotient
is the smash product of the individual quotients.  Choosing all
collapses as products of the same one-dimensional collapse makes the
identity of maps literal.  Composing with units and multiplying in
$MU$ gives the last assertion.
\end{proof}

\subsection{Coordinate complements in boxes}

We record explicitly both the local duality and its compatibility.
The equality-collapse pairing is a local version of the classical
collapse and stable-duality constructions \cite{Thom1954,Atiyah1961Thom}.
We prove it here for coordinate arrangements, starting with one plane
and passing to unions by excision.  The proof will also identify the
variance needed to assemble maps over the simplices of $P$.
Write $T=[-2,2]^M$.  A \emph{coordinate arrangement} in $T$ is a finite
union of sets
\[
 L_J=\{z\in T:z_i=0\text{ for all }i\in J\},
 \qquad J\subseteq\{1,\ldots,M\}.
\]

For a PL map $f:P\to C(t,n)$ from a finite polyhedron and a coordinate
arrangement $Z$, our goal is to turn $\omega_{\mathbb S}(f,Z)=0$
into a displacement of size less
than $1/2$ avoiding $Z$, provided
$\dim Z\le M-q$ and $\dim P\le2q-5$ for an integer $q\ge1$.
The local complements will be $(q-2)$-connected.  Equality-collapse
duality will identify the stable obstruction, and this connectivity
will let us realize a compatible stable lift by actual maps in
\Cref{lem:stable-avoidance}.

Consider finitely many pairs of full-dimensional closed axis-parallel
boxes
\begin{equation}\label{eq:box-pairs}
 B_\alpha\subset\operatorname{int}B_\alpha^+
 \subset(-2,2)^M.
\end{equation}
Assume that no coordinate endpoint is zero, and that corresponding
coordinate intervals of $B_\alpha$ and $B_\alpha^+$ either both contain
zero or both miss zero.  In particular they meet precisely the same
coordinate subspaces.  Whenever the boxes are indexed by a partial
order, we also require both families to be increasing in that order.

Choose $\epsilon>0$ smaller than the absolute value of every endpoint
of every box.  If $Z$ is a nonempty arrangement other than $T$, write
$Z=\bigcup_{J\in\mathcal J}L_J$ with every $J$ nonempty, and put
\[
 d_Z(y)=\min_{J\in\mathcal J}\max_{i\in J}|y_i|,
 \qquad U_\epsilon(Z)=\{y\in T:d_Z(y)<\epsilon\},
 \qquad V_\alpha(Z)=B_\alpha\setminus U_\epsilon(Z).
\]
Here $d_Z$ is exactly the maximum-norm distance to $Z$: setting the
specified coordinates to zero stays inside $T$.  For the empty and full
arrangements put $U_\epsilon(\varnothing)=\varnothing$ and
$U_\epsilon(T)=T$.

\begin{lemma}[Thickened complements]\label{lem:coordinate-thickening}
For sufficiently small $\epsilon$ as above, the inclusion
\[
 V_\alpha(Z)\longrightarrow B_\alpha\setminus Z
\]
is a strong deformation retract equivalence, simultaneously for all the
boxes.  The homotopy fixes $V_\alpha(Z)$ pointwise, preserves each box,
and preserves the set of zero coordinates.  If
$\dim Z\le M-q$ for an integer $q\ge1$, then a nonempty such complement is
$(q-2)$-connected.
\end{lemma}

\begin{proof}
The empty and full arrangements are immediate.  In the remaining case,
for $y\notin Z$ set
\[
 \lambda(y)=\max\{1,\epsilon/d_Z(y)\},
 \qquad \lambda_u(y)=1+u(\lambda(y)-1),\quad 0\le u\le1.
\]
Define a coordinatewise homotopy by
\begin{equation}\label{eq:coordinate-expansion}
 H_u(y)_i=\operatorname{sgn}(y_i)
 \min\{\lambda_u(y)|y_i|,\max\{|y_i|,\epsilon\}\},
\end{equation}
with value zero when $y_i=0$.  It is continuous on $T\setminus Z$:
$\lambda_u$ is locally bounded there, which also gives continuity at
zero coordinates.  Coordinates of magnitude at least $\epsilon$ are
fixed.  Each smaller coordinate is expanded, without changing sign,
by at most enough to reach magnitude $\epsilon$.

Every coordinate interval of a box either contains
$[-\epsilon,\epsilon]$ or lies outside that interval.  Thus
\Cref{eq:coordinate-expansion} preserves even asymmetric boxes.
Coordinate magnitudes never decrease, and zero coordinates remain zero,
so the homotopy stays in the complement and preserves
$V_\alpha(Z)$ at every time.  In fact it fixes that subspace pointwise,
since $d_Z(y)\ge\epsilon$ implies $\lambda(y)=1$.
If $d_Z(y)<\epsilon$, each set $J$ contains an index with
$|y_i|\ge d_Z(y)$, and that coordinate has magnitude at least
$\epsilon$ at time one.  Hence $d_Z(H_1(y))\ge\epsilon$.
This proves the retraction.  The formula is the same for every box.
In particular it applies to any compact family in the unthickened
complement, with no need to move arbitrarily small nonzero coordinates
discontinuously to a fixed radius.

For connectivity, first move a compact family from the closed box into
its interior without changing its zero coordinates: shrink intervals
containing zero by a fixed factor in $(0,1)$, and move the other intervals toward an
interior point of the same sign.  A map from $S^i$ into the complement,
$i\le q-2$, can then be approximated by a PL map in the interior and
extended over a triangulated $D^{i+1}$ into the convex box.  A sufficiently
small generic perturbation of its vertices misses all the excluded
subspaces, since
\[
 (i+1)+(M-q)<M.
\]
The perturbed boundary map is homotopic to the original within the
complement.  This proves the claimed connectivity, including path
connectedness in the applicable range.  Existence of a point in the
complement when $q\ge1$ follows from the same dimension observation.
The retraction transfers connectivity to $V_\alpha(Z)$.
\end{proof}

The local problem is to replace a map into $B_\alpha$ by one with
values in $V_\alpha(Z)$.  We first seek a lift after stabilization.
The cofiber of the complement inclusion records its obstruction;
the corresponding point test retains the parameters of $Z$ inside
$B_\alpha^+$.  Define finite spectra and finite pointed test spaces by
\begin{align}
 \mathcal C_\alpha(Z)
 &=\operatorname{cofib}\left(
 \Sigma^\infty V_\alpha(Z)_+\longrightarrow
 \Sigma^\infty(B_\alpha)_+\right),\label{eq:local-cofiber}\\
 T_\alpha(Z)&=Z/(Z\setminus\operatorname{int}B_\alpha^+).
 \label{eq:local-test-space}
\end{align}
The letter $T_\alpha$ in \Cref{eq:local-test-space} denotes a pointed
quotient, rather than the parameter cube $T$.

\begin{lemma}[Compatible local duality]\label{lem:box-duality}
For the finite box families in \Cref{eq:box-pairs}, one sufficiently
small collapse support gives equivalences
\begin{equation}\label{eq:box-duality}
 \mathcal C_\alpha(Z)\simeq
 F\left(\Sigma^\infty T_\alpha(Z),\Sigma^M\mathbb S\right).
\end{equation}
They are induced by collapse around equality and are natural both under
box inclusions and under inclusions of coordinate arrangements.  In the
box variable the cofibers are covariant and the test spaces are
contravariant, by quotient collapse.  In the arrangement variable,
the cofibers are contravariant and the test spaces are covariant.
\end{lemma}

\begin{proof}
Choose one $\rho>0$ with
\begin{equation}\label{eq:common-collapse-radius}
 \rho<\epsilon,\qquad
 \rho<\min_\alpha
 \operatorname{dist}_\infty
 \bigl(B_\alpha,T\setminus\operatorname{int}B_\alpha^+\bigr).
\end{equation}
All minima are positive, by finiteness and strict containment.  The
formula $(y,z)\mapsto\kappa_\rho(y-z)$ descends to a pairing
\begin{equation}\label{eq:equality-pairing}
 \mathcal C_\alpha(Z)\wedge\Sigma^\infty T_\alpha(Z)
 \longrightarrow\Sigma^M\mathbb S.
\end{equation}
Indeed, if $y\in V_\alpha(Z)$ and $z\in Z$, their distance is at least
$\epsilon$.  If $z\notin\operatorname{int}B_\alpha^+$ and
$y\in B_\alpha$, the second inequality in
\Cref{eq:common-collapse-radius} applies.  These are exactly the two
collapsed subspaces.  Since $V_\alpha(Z)$ is a subpolyhedron, its
ordinary quotient by inclusion in $B_\alpha$ represents the cofiber
in \Cref{eq:local-cofiber}.

First take $Z=L_J$, with $r=|J|$ and $b=M-r$.  If it misses
$B_\alpha$, some required coordinate is bounded away from zero by more
than $\epsilon$, so $V_\alpha(Z)=B_\alpha$.  It also misses
$B_\alpha^+$, and both sides of \Cref{eq:box-duality} are zero.
Otherwise split the box into its normal and tangential factors,
$B_\alpha=B_N\times B_T$.  Then
\[
 B_\alpha/V_\alpha(L_J)
 \cong
 \left(B_N/(B_N\setminus(-\epsilon,\epsilon)^r)\right)
 \wedge(B_T)_+\simeq S^r,
 \qquad T_\alpha(L_J)\cong S^b.
\]
The conventions include $r=0$ and $b=0$.  Fixing an interior point
$w\in B_T$ identifies the pairing with the collapse of
$(u,z_T)\mapsto(u,w-z_T)$, up to a coordinate permutation.  It has
degree $\pm1$: the origin has one preimage, and its local linear map is
$\operatorname{diag}(I_r,-I_b)$.  Its adjoint is therefore an
equivalence $S^r\simeq F(S^b,\Sigma^M\mathbb S)$.

Here is the excision argument for unions.  For coordinate planes,
\[
 U_\epsilon(L_J)\cap U_\epsilon(L_K)
 =U_\epsilon(L_{J\cup K}).
\]
Distributing finite unions gives, for arrangements $Z,Z'$, the
identities
\begin{align*}
 V_\alpha(Z\cup Z')&=V_\alpha(Z)\cap V_\alpha(Z'),\\
 V_\alpha(Z\cap Z')&=V_\alpha(Z)\cup V_\alpha(Z').
\end{align*}
A rectangular subdivision followed by triangulation makes all these
closed complements subpolyhedra.  Their union square is a homotopy
pushout.  Applying suspension spectra and cofibers into the constant
spectrum $\Sigma^\infty(B_\alpha)_+$ gives the homotopy pullback
description
\begin{equation}\label{eq:cofiber-arrangement-excision}
 \mathcal C_\alpha(Z\cup Z')\simeq
 \mathcal C_\alpha(Z)
 \mathop{\times}\limits^{h}_{\mathcal C_\alpha(Z\cap Z')}
 \mathcal C_\alpha(Z').
\end{equation}
We used here that a pushout of spectra is also a pullback.
On the test-space side there is a pushout of pointed polyhedra along
cofibrations
\[
 T_\alpha(Z\cup Z')\cong
 T_\alpha(Z)\cup_{T_\alpha(Z\cap Z')}T_\alpha(Z').
\]
Its function spectra give the same pullback as in
\Cref{eq:cofiber-arrangement-excision}, and the pairings commute with
these maps.  Induction on the number of planes proves
\Cref{eq:box-duality}: the intersection of one plane with a union of
fewer planes is again a union of no more than that many coordinate
planes.  The empty arrangement contributes zero throughout.

Finally, for an arrangement inclusion $Z\subseteq Z'$, the maps are
\[
 \mathcal C_\alpha(Z')\longrightarrow\mathcal C_\alpha(Z),
 \qquad T_\alpha(Z)\longrightarrow T_\alpha(Z').
\]
For a box inclusion indexed by $\alpha\le\beta$, the maps are
\[
 \mathcal C_\alpha(Z)\longrightarrow\mathcal C_\beta(Z),
 \qquad T_\beta(Z)\longrightarrow T_\alpha(Z),
\]
the second being quotient collapse.  On representatives, each required
compatibility says that the pairing is $\kappa_\rho(y-z)$.  If the
smaller test quotient additionally collapses $z$, the pairing with
$y\in B_\alpha$ was already the basepoint by
\Cref{eq:common-collapse-radius}.  Thus the adjoints are natural
equivalences of diagrams, not just objectwise identifications.
\end{proof}

\subsection{From a stable lift to a map of complements}

The local duality is now available as an equivalence of diagrams.
We use a single global nullhomotopy to lift all simplex maps
simultaneously, and then use connectivity to realize the lift by
actual maps into the complements.

\begin{lemma}[Stable avoidance]\label{lem:stable-avoidance}
Let $P$ be a finite polyhedron and $f:P\to C(t,n)$ a PL map.
Let $Z\subseteq T(t,n)$ be a coordinate arrangement, and let $q\ge1$
be an integer.  Suppose that
\[
 \omega_{\mathbb S}(f,Z)=0,\qquad
 \dim Z\le M-q,\qquad \dim P\le2q-5.
\]
There is a continuous map
\[
 g:P\longrightarrow(-2,2)^M\setminus Z
 \quad\text{such that}\quad
 \|g(x)-f(x)\|_\infty<\tfrac12\quad(x\in P).
\]
\end{lemma}

\begin{proof}
The empty domain or empty arrangement needs no argument.  Otherwise the
dimension hypotheses imply $q\ge3$.

\emph{Choice of boxes.}
Subdivide $P$ so finely that $f(\sigma)$ has maximum-norm diameter less
than $c$ for each simplex.  Thicken its coordinate bounding box by a
radius between $2c$ and $3c$, increasing strictly with the dimension of
$\sigma$.  A face box is then contained in the interior of its coface
box.  Choose the finitely many radii generically so that no endpoint is
zero.  The resulting boxes $B_\sigma$ have the properties
\begin{equation}\label{eq:avoidance-box-control}
 \operatorname{dist}_\infty
 \bigl(f(\sigma),T\setminus\operatorname{int}B_\sigma\bigr)>c,
 \qquad
 \sup_{x\in\sigma,\,y\in B_\sigma}\|y-f(x)\|_\infty<\tfrac12.
\end{equation}
They all lie in $(-2,2)^M$; indeed the latter supremum is less than
$4c$.  Enlarge them by one sufficiently small uniform amount to obtain
$B_\sigma^+$, preserving the coordinate incidences and all inclusions.
Choose a common $\epsilon$ and then a common collapse radius
$\rho<c$ as in \Cref{lem:coordinate-thickening,lem:box-duality}.
Write $V_\sigma=V_\sigma(Z)$ and $T_\sigma=T_\sigma(Z)$.

\emph{The stable lift as a map of diagrams.}
Index covariant diagrams by the nonempty simplices of $P$, with maps
from faces to cofaces.  Set
\[
 \mathcal D_\sigma=\Sigma^\infty\sigma_+,
 \qquad \mathcal B_\sigma=\Sigma^\infty(B_\sigma)_+,
 \qquad \mathcal V_\sigma=\Sigma^\infty(V_\sigma)_+.
\]
The maps $f|_\sigma$ define a natural map
$\mathcal D\to\mathcal B$.  Its obstruction to lifting through
$\mathcal V\to\mathcal B$ is the composite to the cofiber diagram
$\mathcal C$.

For clarity, the source cell diagram has latching maps
\begin{equation}\label{eq:cell-latching}
 \underset{\tau\subsetneq\sigma}{\operatorname{colim}}\tau_+
 =\partial\sigma_+\longrightarrow\sigma_+.
\end{equation}
They are cofibrations; at a vertex the latching object is the basepoint.
Consequently this is a cofibrant diagram, and compatible mapping out
of its suspension spectrum computes the derived mapping spectrum.
This fact can also be seen directly: add the simplices in increasing
dimension.  At each step a compatible map is an extension over
$\sigma$ of the already specified map on $\partial\sigma$.
Restriction along \Cref{eq:cell-latching} is a mapping-space fibration
after fibrant replacement.  Induction therefore makes compatible
mapping invariant under objectwise equivalences and preserves the
fiber sequences of mapping spectra used here.

The duality in \Cref{lem:box-duality} identifies the obstruction with
the adjoint of the equality test on $\sigma_+\wedge T_\sigma$.
There is a well-defined based map
\begin{equation}\label{eq:local-to-global-test}
 \sigma_+\wedge T_\sigma\longrightarrow Q_c(f,Z),
 \qquad [x,z]\longmapsto[x,z].
\end{equation}
Indeed, if $z$ lies outside $\operatorname{int}B_\sigma^+$, then
\Cref{eq:avoidance-box-control} puts $(x,z)$ in the relative set of
the global test.  For a face $\sigma\subset\tau$, compare the two
maps on $\sigma_+\wedge T_\tau$: including $\sigma$ in $\tau$ or
collapsing $T_\tau$ to $T_\sigma$ gives the same based map to
$Q_c(f,Z)$.
Every additional collapse already lands in the global relative set by
\Cref{eq:avoidance-box-control}.  After adjunction these equalities
give the required natural transformation.  Pulling back one stable nullhomotopy of
$\varphi_\rho(f,Z)$ thus gives a compatible nullhomotopy of the
obstruction diagram.  This uses the assumed relative point-test
vanishing, rather than separately chosen local nullhomotopies.

The category of spectrum diagrams is stable.  Mapping out of
$\mathcal D$ into the objectwise cofiber sequence
$\mathcal V\to\mathcal B\to\mathcal C$ now gives a stable lift
$\mathcal D\to\mathcal V$, together with a compatible homotopy of its
composite to the map given by $f$.

\emph{Augmentation and desuspension.}
For an unbased space $V$, the augmentation
$\Sigma^\infty V_+\to\mathbb S$ sends every point to the nonbasepoint
of $S^0$.  A map of spaces into $V$ therefore has augmentation equal
to the unit, whereas an arbitrary stable map need not.
Each $B_\sigma$ is contractible, so its augmentation is an equivalence.
After adjunction, augmenting the compatible homotopy in the box
diagram gives a homotopy from the stable lift's augmentation to the
constant unit map.
We retain the lift together with this homotopy by setting
\begin{equation}\label{eq:augmentation-fiber}
 E_\sigma=\operatorname{hofib}_{1}\left(
 \Omega^\infty\Sigma^\infty(V_\sigma)_+
 \longrightarrow\Omega^\infty\mathbb S\right),
\end{equation}
where $1$ is the unit point.  Naturality of the augmentations makes
these data a single derived map from the cell diagram $\sigma$
into the homotopy-fiber diagram $E$.  We keep the lift and its
augmentation homotopy together when choosing functorial models.

There is a natural comparison $j_\sigma:V_\sigma\to E_\sigma$.
Its augmentation is literally the unit, so the path in the homotopy
fiber can be taken constant.  To compute its connectivity, choose a
point $v_\sigma\in V_\sigma$.  The resulting splitting of the
augmentation identifies $E_\sigma$ with
$\Omega^\infty\Sigma^\infty(V_\sigma,v_\sigma)$ and identifies
$j_\sigma$ with stabilization.  The choices are used only for this
objectwise connectivity calculation; the comparison itself was
defined naturally.  By \Cref{lem:coordinate-thickening},
$V_\sigma$ is $(q-2)$-connected.  Freudenthal therefore makes
$j_\sigma$ at least $(2q-3)$-connected
\cite[Corollary~4.24]{Hatcher2002}.  In particular, the weaker bound
$2q-4$ is available with room for the stated dimension hypothesis.

Here is an explicit way to realize the lift without making compatible
basepoint choices.  Replace $j$ by its natural mapping-path
factorization.  Objectwise this is
\[
 W_\sigma=\{(v,\gamma):v\in V_\sigma,\ 
 \gamma:[0,1]\to E_\sigma,\ \gamma(0)=j_\sigma(v)\},
 \qquad p_\sigma(v,\gamma)=\gamma(1).
\]
The maps $p_\sigma$ are fibrations of connectivity at least $2q-3$.
Projection $W_\sigma\to V_\sigma$ is a natural homotopy equivalence,
with section given by constant paths.  The cofibrant cell diagram
allows the preceding compatible derived map to $E$ to be represented
by an actual map of diagrams, using objectwise fibrant models.

Inductively lift this map to $W$.  At a simplex $\sigma$, the chosen
lifts on proper faces give a map $\partial\sigma\to W_\sigma$ over
the prescribed map $\sigma\to E_\sigma$.  Its extension is possible
because the fibers of $p_\sigma$ are $(2q-4)$-connected and
$\dim\sigma\le2q-5$.  This is the ordinary lifting obstruction for
a disk and its boundary.  Projection back to $V$ gives an actual
compatible family $\sigma\to V_\sigma$.  No dimension of a separate
indexing space is added: the extension domains are exactly the
simplices in \Cref{eq:cell-latching}.

The family glues to a continuous map $g:P\to(-2,2)^M\setminus Z$.
Its image on $\sigma$ lies in $B_\sigma$, so
\Cref{eq:avoidance-box-control} gives the claimed displacement bound.
\end{proof}

\subsection{Restoring cube boundaries}

The next elementary step is stated for compact inputs, since no
polyhedral hypothesis is needed.  Its preservation clause is what
allows successive applications of \Cref{lem:stable-avoidance}.

\begin{lemma}[Boundary restoration]\label{lem:boundary-restoration}
Let $K$ be compact, let $0\le s<n$ be an integer, and suppose that
$f:K\to C(t,n)$ and
$g:K\to(-2,2)^M\setminus Z_s(t,n)$ are continuous and satisfy
$\|g(x)-f(x)\|_\infty<1/2$ for every $x$.
There is a continuous map $G:K\to C(t,n)$ such that at least $s+1$
of its slots belong to $\partial[-1,1]^t$ at every point, and
\[
 f_j(x)\in\partial[-1,1]^t\quad\Longrightarrow\quad G_j(x)=f_j(x)
 \qquad(1\le j\le n).
\]
\end{lemma}

\begin{proof}
If $K$ is empty the assertion is immediate.  Otherwise,
choose a continuous $\chi:[0,1]\to[0,1]$ which is zero on $[0,3/4]$
and equals one at $1$.  Define slotwise
\[
 g'_j(x)=\bigl(1-\chi(\|f_j(x)\|_\infty)\bigr)g_j(x)
       +\chi(\|f_j(x)\|_\infty)f_j(x).
\]
Whenever the cutoff is nonzero, one scalar coordinate of $f_j(x)$ has
magnitude greater than $3/4$.  The corresponding coordinate of
$g_j(x)$ has the same sign and magnitude greater than $1/4$.
The convex interpolation consequently cannot turn that slot into zero.
All other slots are unchanged.  Hence $g'$ still has at least $s+1$
nonzero slots everywhere, and it already agrees with $f$ at each
original boundary slot.

The $(s+1)$-st largest of the finitely many continuous functions
$\|g'_j(x)\|_\infty$ is continuous and positive on $K$.
By compactness choose $0<\delta<1$ smaller than its minimum.  Set
\[
 G_j(x)=\frac{g'_j(x)}{\max\{\delta,\|g'_j(x)\|_\infty\}}.
\]
This is continuous even at zero, belongs to the unit cube, and sends
every slot of radius at least $\delta$ to the boundary.  It fixes
every slot already of radius one.  Both conclusions follow.
\end{proof}
